\chapter{Задача первого исследовательского цикла}

\section{Результат, на который направлена работа}
Подготовить записку с одной основной теоремой, полным доказательством,
проверенным примером и точным сравнением с ближайшей литературой.
Основная теорема должна решать явно сформулированную трудность:
давать новую конструкцию, оценку, класс применимости или препятствие
для существующего метода.

Учебные леммы и воспроизведения известных доказательств включаются в
рабочие записи с указанием источника. Статус основного результата
пересматривается после поиска предшественников и внешнего обсуждения.

Первый цикл занимает 12 учебных недель, то есть 144--180 часов.
Контрольные этапы --- конец недель~2, 4, 8 и~12.
Выбор задачи начинается в первый рабочий день.

\section{Что уже есть в репозитории}
\begin{description}
\item[Достаточный топологический критерий.]
\repo{Mlc/ParameterComponentApproximation.lean}:
буферное условие для убывающих компактных приближений влечёт локальную
связность предела. Для множества Мандельброта соответствующий вход
передаётся явно.
\item[Форма конечного покрытия.]
\repo{Mlc/FiniteComponentCriterion.lean}:
поле \repo{LocalPiece.control} всё ещё содержит
\(\forall N\,\exists L\) и включение в компоненту.
Для него требуется отдельное математическое обоснование.
\item[Доступные конструкции.]
Внешние орбитальные приближения, рациональные прямоугольники, интервальные
оценки и ловушки. Начальные модули ---
\repo{Mlc/CertifiedOrbitApproximation.lean},
\repo{Mlc/CertifiedTrappingRegions.lean},
\repo{Mlc/FiniteTrappingRegions.lean}.
\item[Установленные ограничения моделей.]
\repo{Mlc/GreenSublevelIntersectionCounterexample.lean} содержит
контрпример к прежнему универсальному утверждению.
\repo{Mlc/Literature/ModelAudit.lean} выявляет последствия подмены
конформного модуля гауссовой площадью.
\end{description}
Для каждого используемого результата записать точную декларацию, автора,
дату и источник по истории проекта. Условные входы раскрыть до определений.

\section{Проверки постановки}
До длинного доказательства выполнить четыре действия:
\begin{enumerate}
\item Предъявить объект, удовлетворяющий всем гипотезам, и вычислить его
      входные данные.
\item Проверить, что гипотезы можно установить без предполагаемого вывода.
\item Найти ближайшую известную теорему и проверить возможность прямого
      применения к выбранному классу.
\item Назвать один переход, для которого в проекте действительно
      отсутствует доказательство.
\end{enumerate}
Наличие слова «конечный» в названии структуры проверяется по типам её полей.
Для алгоритма отдельно проверяются представление входа, завершаемость
и смысл выдаваемого ответа.

\section{Подготовка по возникающей зависимости}
Для недостающей леммы записать: место её использования, точную формулировку,
источник и требуемый объём разбора. Сначала попытаться доказать нужный
частный случай. Если требуется крупная внешняя теорема, привести её
формулировку и проверить гипотезы.

Механический ознакомительный проход --- медленное чтение и воспроизведение
обозначений и полных формулировок. Он помогает освоить форму незнакомого
фрагмента. В этом цикле такой проход занимает до 30 минут внутри
читательского бюджета и заканчивается указанием нужной леммы.
Непрерывное чтение Любича после с.~133 временно снимается с расписания.
Михайлов--Пасси, Хэтчер и Милнор подключаются при появлении соответствующих
зависимостей.
